% Part: first-order-logic
% Chapter: sequent-calculus
% Section: proving-things

\documentclass[../../../include/open-logic-section]{subfiles}

\begin{document}

\iftag{FOL}
      {\olfileid{fol}{seq}{pro}}
      {\olfileid{pl}{seq}{pro}}

\olsection{Tuladha \usetoken{P}{derivation}}

\begin{ex}
Wenehana $\Log{LK}$-!!{derivation} kanggo sekuen $!A \land !B \Sequent !A$.

Kita miwiti kanthi nulis sekuen-pungkasan kang dikarepake ing sisih
paling ngisor !!{derivation}.
\begin{prooftree}
\AxiomC{}
\UnaryInf$!A\land !B \fCenter !A$
\end{prooftree}
Sabanjure, kita kudu nemtokake jinis inferensi kang bisa nduweni
sekuen ngisor kanthi wujud iki. Iki bisa wae aturan struktural, nanging
becike kita miwiti kanthi nggoleki aturan logis. Konektif logis
siji-sijine kang muncul ing sekuen ngisor yaiku $\land$, mula kita
nggoleki aturan $\land$, lan amarga simbol $\land$ muncul ing
anteseden, kita nggatekake aturan \LeftR{\land}.
\begin{prooftree}
\AxiomC{}
\RightLabel{\LeftR{\land}}
\UnaryInf$!A\land !B \fCenter !A$
\end{prooftree}
Ana rong pilihan kanggo sekuen ndhuwur ing inferensi
\LeftR{\land}: sekuen ndhuwur iku bisa $!A \Sequent !A$, utawa $!B
\Sequent !A$. Temenan, $!A \Sequent !A$ iku sekuen wiwitan (iki
apik), dene $!B \Sequent !A$ ora mesthi bisa diderivasi. Kita ngisi
sekuen ndhuwur:
\begin{prooftree}
\Axiom$!A \fCenter !A$
\RightLabel{\LeftR{\land}}
\UnaryInf$!A\land !B \fCenter !A$
\end{prooftree}
Saiki kita duwe $\Log{LK}$-!!{derivation} kang bener kanggo sekuen $!A
\land !B \Sequent !A$.
\end{ex}

\begin{ex}
Wenehana $\Log{LK}$-!!{derivation} kanggo sekuen $\lnot !A \lor !B
\Sequent !A \lif !B$.

Miwitia kanthi nulis sekuen-pungkasan kang dikarepake ing sisih paling
ngisor !!{derivation}.
\begin{prooftree}
\AxiomC{}
\UnaryInf$\lnot !A \lor !B \fCenter !A \lif !B$
\end{prooftree}
Kanggo nemokake aturan logis kang bisa ngasilake sekuen-pungkasan iki,
kita ndeleng konektif logis ing jerone: $\lnot$, $\lor$, lan $\lif$.
Saiki kita mung nggatekake $\lor$ lan $\lif$ amarga kalorone iku
!!{main operator}s saka !!{sentence}s ing sekuen-pungkasan, dene
$\lnot$ ana ing jangkauan konektif liyane, mula bakal ditangani
mengko. Dadi, pilihan aturan logis kanggo inferensi pungkasan yaiku
aturan \LeftR{\lor} lan aturan \RightR{\lif}. Sejatine kita bisa milih
salah siji, nanging ayo milih aturan \RightR{\lif} (saora-orane amarga
aturan iki ngidini kita nundha mecah dadi rong cabang). Miturut wujud
inferensi \RightR{\lif} kang bisa ngasilake sekuen ngisor iki,
wujude kudu kaya mangkene:
\begin{prooftree}
\AxiomC{}
\UnaryInf$ !A, \lnot !A \lor !B \fCenter !B $
\RightLabel{\RightR{\lif}} \UnaryInf$ \lnot !A \lor !B \fCenter !A \lif !B $
\end{prooftree}

Yen $\lnot !A \lor !B$ dipindhah menyang pucuk njaba anteseden, kita
bisa ngetrapake aturan \LeftR{\lor}. Miturut skemane, inferensi iki
kudu pecah dadi rong sekuen ndhuwur kaya mangkene:
\begin{prooftree}
\AxiomC{}
\UnaryInf$\lnot !A, !A \fCenter !B$
\AxiomC{}
\UnaryInf$!B, !A \fCenter !B$
\RightLabel{\LeftR{\lor}}
\BinaryInf$ \lnot !A \lor !B, !A \fCenter !B $
\RightLabel{\RightR{\Exchange}}
\UnaryInf$ !A, \lnot !A \lor !B \fCenter !B $
\RightLabel{\RightR{\lif}}
\UnaryInf$ \lnot !A \lor !B \fCenter !A \lif !B $
\end{prooftree}
Elinga yen kita ngupaya nglacak munggah nganti tekan sekuen wiwitan;
kayane kita wis cedhak!{} Cabang tengen mung kurang siji pelemahan lan
siji pertukaran saka sekuen wiwitan, banjur cabang iku rampung:
\begin{prooftree}
\AxiomC{}
\UnaryInf$\lnot !A, !A \fCenter !B$
\Axiom$!B \fCenter !B$
\RightLabel{\LeftR{\Weakening}}
\UnaryInf$!A, !B \fCenter !B$
\RightLabel{\LeftR{\Exchange}}
\UnaryInf$!B, !A \fCenter !B$
\RightLabel{\LeftR{\lor}}
\BinaryInf$\lnot !A \lor !B, !A \fCenter !B $
\RightLabel{\RightR{\Exchange}}
\UnaryInf$ !A, \lnot !A \lor !B \fCenter !B $
\RightLabel{\RightR{\lif}}
\UnaryInf$ \lnot !A \lor !B \fCenter !A \lif !B $
\end{prooftree}

Saiki delengen cabang kiwa. Konektif logis siji-sijine ing sembarang
!!{sentence} yaiku simbol $\lnot$ ing !!{sentence}s anteseden, mula
kita nggatekake sawijining kedadeyan aturan \LeftR{\lnot}.
\begin{prooftree}
\AxiomC{}
\UnaryInf$ !A \fCenter !B, !A$
\RightLabel{\LeftR{\lnot}}
\UnaryInf$\lnot !A, !A \fCenter !B$
\Axiom$!B \fCenter !B$
\RightLabel{\LeftR{\Weakening}}
\UnaryInf$!A, !B \fCenter !B$
\RightLabel{\LeftR{\Exchange}}
\UnaryInf$!B, !A \fCenter !B$
\RightLabel{\LeftR{\lor}}
\BinaryInf$\lnot !A \lor !B, !A \fCenter !B $
\RightLabel{\RightR{\Exchange}}
\UnaryInf$ !A, \lnot !A \lor !B \fCenter !B $
\RightLabel{\RightR{\lif}}
\UnaryInf$ \lnot !A \lor !B \fCenter !A \lif !B $
\end{prooftree}
Kaya nalika ngrampungake cabang tengen, kita mung kurang siji
pelemahan lan siji pertukaran kanggo ngrampungake cabang kiwa iki.
\begin{prooftree}
\Axiom$!A \fCenter !A$
\RightLabel{\RightR{\Weakening}}
\UnaryInf$ !A \fCenter !A, !B$
\RightLabel{\RightR{\Exchange}}
\UnaryInf$ !A \fCenter !B, !A$
\RightLabel{\LeftR{\lnot}}
\UnaryInf$\lnot !A, !A \fCenter !B$
\Axiom$!B \fCenter !B$
\RightLabel{\LeftR{\Weakening}}
\UnaryInf$!A, !B \fCenter !B$
\RightLabel{\LeftR{\Exchange}}
\UnaryInf$!B, !A \fCenter !B$
\RightLabel{\LeftR{\lor}}
\BinaryInf$\lnot !A \lor !B, !A \fCenter !B $
\RightLabel{\RightR{\Exchange}}
\UnaryInf$ !A, \lnot !A \lor !B \fCenter !B $
\RightLabel{\RightR{\lif}}
\UnaryInf$ \lnot !A \lor !B \fCenter !A \lif !B $
\end{prooftree}
\end{ex}

\begin{ex}
Wenehana $\Log{LK}$-!!{derivation} kanggo sekuen $\lnot !A \lor \lnot !B
\Sequent \lnot (!A \land !B)$

Kanthi teknik ing ndhuwur, kita miwiti kanthi nulis sekuen-pungkasan
kang dikarepake ing sisih paling ngisor.
\begin{prooftree}
\AxiomC{}
\UnaryInf$ \lnot !A \lor \lnot !B \fCenter \lnot (!A \land !B) $
\end{prooftree}
Konektif utama saka !!{sentence}s ing sekuen-pungkasan yaiku simbol
$\lor$ lan simbol $\lnot$. Ing kene aturan \LeftR{\lor} utawa aturan
\RightR{\lnot} bisa ditrapake, nanging kita miwiti nganggo aturan
\RightR{\lnot} amarga sauntara ora perlu mecah dadi rong cabang:
\begin{prooftree}
\AxiomC{}
\UnaryInf$!A \land !B, \lnot !A \lor \lnot !B \fCenter $
\RightLabel{\RightR{\lnot}}
\UnaryInf$\lnot !A \lor \lnot !B \fCenter \lnot (!A \land !B)$
\end{prooftree}
Saiki kita bisa milih aturan \LeftR{\land} utawa \LeftR{\lor}. Ayo
deleng apa kang kelakon yen aturan \LeftR{\land} ditrapake: kita bisa
miwiti saka sekuen $!A, \lnot !A \lor \lnot !B \Sequent \quad$ utawa
sekuen $!B, \lnot !A \lor \lnot !B \Sequent \quad$. Amarga
!!{derivation} iku simetris tumrap $!A$ lan $!B$, ayo milih kang
kapisan. Cathetan koreksi sumber OLPL-004: loro formula ing prosa
Inggris ilang negasi ing disjung kapindho; negasi kasebut dibalekake
supaya cocog karo sekuen-pungkasan lan wit kang langsung nututi.
\begin{prooftree}
\AxiomC{}
\UnaryInf$!A, \lnot !A \lor \lnot !B \fCenter $
\RightLabel{\LeftR{\land}}
\UnaryInf$!A \land !B, \lnot !A \lor \lnot !B \fCenter $
\RightLabel{\RightR{\lnot}}
\UnaryInf$\lnot !A \lor \lnot !B \fCenter \lnot (!A \land !B)$
\end{prooftree}
Nalika terus ngisi !!{derivation}, kita weruh yen ana masalah:
\begin{prooftree}
\Axiom$!A \fCenter !A$
\RightLabel{\LeftR{\lnot}}
\UnaryInf$ \lnot !A, !A \fCenter$
\AxiomC{}
\RightLabel{?}
\UnaryInf$!A \fCenter !B$
\RightLabel{\LeftR{\lnot}}
\UnaryInf$ \lnot !B, !A \fCenter$
\RightLabel{\LeftR{\lor}}
\BinaryInf$\lnot !A \lor \lnot !B, !A \fCenter $
\RightLabel{\LeftR{\Exchange}}
\UnaryInf$!A, \lnot !A \lor \lnot !B \fCenter $
\RightLabel{\LeftR{\land}}
\UnaryInf$!A \land !B, \lnot !A \lor \lnot !B \fCenter $
\RightLabel{\RightR{\lnot}}
\UnaryInf$\lnot !A \lor \lnot !B \fCenter \lnot (!A \land !B)$
\end{prooftree}
Pucuk cabang tengen ora bisa dikurangi maneh lan ora bisa digawa
menyang sekuen wiwitan nganggo inferensi struktural, mula iki dudu
dalan kang bener. Dadi, penerapan aturan \LeftR{\land} ing ndhuwur
temenan kleru. Yen kita bali menyang langkah sadurunge lan ngetrapake
aturan \LeftR{\lor}, asile
\begin{prooftree}
\AxiomC{}
\UnaryInf$\lnot !A, !A \land !B \fCenter $

\AxiomC{}
\UnaryInf$\lnot !B, !A \land !B \fCenter $

\RightLabel{\LeftR{\lor}}
\BinaryInf$\lnot !A \lor \lnot !B, !A \land !B \fCenter $
\RightLabel{\LeftR{\Exchange}}
\UnaryInf$!A \land !B, \lnot !A \lor \lnot !B \fCenter $
\RightLabel{\RightR{\lnot}}
\UnaryInf$\lnot !A \lor \lnot !B \fCenter \lnot (!A \land !B)$
\end{prooftree}
Kanthi ngrampungake saben cabang kaya sadurunge, kita entuk
\begin{prooftree}
\Axiom$ !A \fCenter!A$
\RightLabel{\LeftR{\land}}
\UnaryInf$!A \land !B \fCenter !A$
\RightLabel{\LeftR{\lnot}}
\UnaryInf$\lnot !A, !A \land !B \fCenter $

\Axiom$ !B \fCenter !B$
\RightLabel{\LeftR{\land}}
\UnaryInf$!A \land !B \fCenter !B$
\RightLabel{\LeftR{\lnot}}
\UnaryInf$\lnot !B, !A \land !B \fCenter $

\RightLabel{\LeftR{\lor}}
\BinaryInf$\lnot !A \lor \lnot !B, !A \land !B \fCenter $
\RightLabel{\LeftR{\Exchange}}
\UnaryInf$!A \land !B, \lnot !A \lor \lnot !B \fCenter $
\RightLabel{\RightR{\lnot}}
\UnaryInf$\lnot !A \lor \lnot !B \fCenter \lnot (!A \land !B)$
\end{prooftree}
(Ing langkah iki, aturan $\land$ bisa ditrapake luwih ngisor tinimbang
aturan $\lnot$ lan kita tetep ngasilake !!{derivation} kang bener).
\end{ex}

\begin{ex}
Nganti saiki kita durung nggunakake aturan kontraksi, nanging aturan
iki kadhangkala dibutuhake. Iki tuladha nalika kedadeyan mangkono.
Upamana kita arep mbuktekake $\quad \Sequent !A \lor \lnot !A$.
Ngetrapake $\RightR{\lor}$ kanthi mbalikke bakal menehi salah siji saka
rong !!{derivation}s iki:
\begin{prooftree}
\AxiomC{}
\UnaryInf$ \fCenter !A$
\RightLabel{\RightR{\lor}}
\UnaryInf$ \fCenter !A \lor \lnot !A$
\DisplayProof\qquad\bottomAlignProof
\AxiomC{}
\UnaryInf$!A \fCenter $
\RightLabel{\RightR{\lnot}}
\UnaryInf$ \fCenter \lnot !A$
\RightLabel{\RightR{\lor}}
\UnaryInf$ \fCenter !A \lor \lnot !A$
\end{prooftree}
Mesthi wae, loro-lorone ora mungkasi ing sekuen wiwitan. Carane yaiku
ngelingi yen aturan kontraksi ngidini kita nggabungake rong salinan
saka !!a{sentence} dadi siji---lan nalika nggoleki bukti, yaiku lumaku
saka ngisor menyang ndhuwur, kita bisa nyimpen salinan $!A \lor \lnot
!A$ ing premis, upamane,
\begin{prooftree}
\AxiomC{}
\UnaryInf$ \fCenter !A \lor \lnot !A, !A$
\RightLabel{\RightR{\lor}}
\UnaryInf$ \fCenter !A \lor \lnot !A, !A \lor \lnot !A$
\RightLabel{\RightR{\Contraction}}
\UnaryInf$ \fCenter !A \lor \lnot !A$
\end{prooftree}
Saiki kita bisa ngetrapake $\RightR{\lor}$ kaping pindho lan uga
entuk~$\lnot !A$, kang ngasilake !!{derivation} jangkep.
\begin{prooftree}
\Axiom$!A \fCenter !A$
\RightLabel{\RightR{\lnot}}
\UnaryInf$\fCenter !A, \lnot !A$
\RightLabel{\RightR{\lor}}
\UnaryInf$\fCenter !A, !A \lor \lnot !A$
\RightLabel{\RightR{\Exchange}}
\UnaryInf$ \fCenter !A \lor \lnot !A, !A$
\RightLabel{\RightR{\lor}}
\UnaryInf$ \fCenter !A \lor \lnot !A, !A \lor \lnot !A$
\RightLabel{\RightR{\Contraction}}
\UnaryInf$ \fCenter !A \lor \lnot !A$
\end{prooftree}
\end{ex}

\begin{prob}
Wenehana !!{derivation}s kanggo sekuen-sekuen iki:
\begin{enumerate}
\item $!A \land (!B \land !C) \Sequent (!A \land !B) \land !C$.
\item $!A \lor (!B \lor !C) \Sequent (!A \lor !B) \lor !C$.
\item $!A \lif (!B \lif !C) \Sequent !B \lif (!A \lif !C)$.
\item $!A \Sequent \lnot\lnot !A$.
\end{enumerate}
\end{prob}

\begin{prob}
Wenehana !!{derivation}s kanggo sekuen-sekuen iki:
\begin{enumerate}
\item $(!A \lor !B) \lif !C \Sequent !A \lif !C$.
\item $(!A \lif !C) \land (!B \lif !C) \Sequent (!A \lor !B) \lif !C$.
\item $\Sequent \lnot(!A \land \lnot !A)$.
\item $!B \lif !A \Sequent \lnot !A \lif \lnot !B$.
\item $\Sequent (!A \lif \lnot !A) \lif \lnot !A$.
\item $\Sequent \lnot(!A \lif !B) \lif \lnot !B$.
\item $!A \lif !C \Sequent \lnot (!A \land \lnot !C)$.
\item $!A \land \lnot !C \Sequent \lnot (!A \lif !C)$.
\item $!A \lor !B, \lnot !B \Sequent !A$.
\item $\lnot !A \lor \lnot !B \Sequent \lnot(!A \land !B)$.
\item $\Sequent (\lnot !A \land \lnot !B) \lif\lnot(!A \lor !B)$.
\item $\Sequent \lnot(!A \lor !B) \lif (\lnot !A \land \lnot !B)$.
\end{enumerate}
\end{prob}

\begin{prob}
Wenehana !!{derivation}s kanggo sekuen-sekuen iki:
\begin{enumerate}
\item $\lnot(!A \lif !B) \Sequent !A$.
\item $\lnot(!A \land !B) \Sequent \lnot !A \lor \lnot !B$.
\item $!A \lif !B \Sequent \lnot !A \lor !B$.
\item $\Sequent \lnot \lnot !A \lif !A$.
\item $!A \lif !B, \lnot !A \lif !B \Sequent !B$.
\item $(!A \land !B) \lif !C \Sequent (!A \lif !C) \lor (!B \lif !C)$.
\item $(!A \lif !B) \lif !A \Sequent !A$.
\item $\Sequent (!A \lif !B) \lor (!B \lif !C)$.
\end{enumerate}
(Kabeh iki mbutuhake aturan~$\RightR{\Contraction}$.)
\end{prob}
\end{document}
